%% ATS Resume Template ("ats-resume")
%% Dense, single-column, plain-LaTeX resume (no moderncv). Modeled on a
%% reference resume that packs far more content per page than moderncv's
%% banking style by eliminating its hint-column date layout and generous
%% inter-entry spacing: entry title + right-aligned date share one line,
%% organization/location sit in italics on the next line, and section
%% headings are a bold caps label immediately followed by a full-width rule.
%%
%% Compile with: lualatex (fontawesome5 icons render cleanly under lualatex;
%% pdflatex frequently fails on modern MiKTeX with font-expansion errors on
%% the same package).
%%
%% Page limit: 1 page (industry resume standard for this framework -- see
%% 05-cv-templates.md "Page Budget - Hard 1-Page Limit"). Every \resumeentry
%% is preceded by \needspace{...} sized for the entry heading plus its first
%% bullet, so an entry can never be orphaned at the bottom of a page with
%% its bullets pushed to the next page.

\documentclass[10pt,letterpaper]{article}

\usepackage[T1]{fontenc}
\usepackage[utf8]{inputenc}
\usepackage{lmodern}
\usepackage{fontawesome5}
\usepackage{xcolor}
% top=0in: kept at 0in even though the bar no longer consumes normal text
% flow (see eso-pic below) -- content still needs to start flush with no
% residual margin, since the manual \vspace after the bar (in the TOP ACCENT
% BAR CLEARANCE block) is what now creates the visual gap before the name.
\usepackage[left=0.65in,right=0.65in,top=0in,bottom=0.5in]{geometry}
% eso-pic draws the accent bar as a page-background layer, independent of
% geometry's margins -- needed so the bar can bleed all the way to the
% physical left/right page edges (2026-08-03, per user feedback that the
% bar needs to "continue on the side" -- the prior \makebox[\linewidth]
% version was capped at the text width, i.e. inside the left/right margins).
\usepackage{eso-pic}
\usepackage{enumitem}
\usepackage{etoolbox}
\usepackage{needspace}
\usepackage{hyperref}

% Accent color applied to name, section headings, and section rules only --
% everything else (contact line, body text, entry titles) stays black for
% ATS-safety and print legibility.
\definecolor{cvblue}{HTML}{1F4E79}

% Top accent bar, full page-width bleed (edge to edge, not capped at the
% text margins). \AtPageUpperLeft places content with its reference point at
% the page's physical top-left corner; \rule's optional negative lift
% (-5pt) then lowers the bar by its own height so it hangs visibly *below*
% that corner instead of extending upward off the physical page (a rule's
% height normally extends upward from its reference point, which would put
% an unlowered bar entirely off-canvas here). 5pt thick (up from the
% text-width version's 3pt) per user request for a "slightly wider" bar.
\AddToShipoutPictureBG{%
  \AtPageUpperLeft{{\color{cvblue}\rule[-5pt]{\paperwidth}{5pt}}}%
}

\hypersetup{
    colorlinks=true,
    linkcolor=cvblue,
    urlcolor=cvblue,
    pdftitle={[YOUR NAME] - Resume},
    pdfpagemode=FullScreen,
}

\pagestyle{empty}
% \raggedbottom (not \flushbottom, the LaTeX default) so the page ends with
% real, unstretched whitespace instead of LaTeX inflating small interline
% gaps to fill the page -- this matters if you ever use \enlargethispage.
\raggedbottom
\setlength{\parindent}{0pt}
\setlength{\parskip}{0pt}
% \topskip (default 10pt) reserves space above the first box/line on every
% page independent of geometry's top margin -- even at top=0in, LaTeX still
% left a ~7pt gap (10pt topskip minus the accent bar's 3pt height) above the
% bar. Set to 0pt so the bar is truly flush with the page edge (2026-08-03).
\setlength{\topskip}{0pt}

% Tight, uniform list spacing -- no \vspace between \item entries anywhere
% in this document (see 05-cv-templates.md "Spacing inside itemize lists").
% topsep MUST stay 0pt: even 1pt here silently adds space before the first
% bullet of every single entry, which reads as "extra space after the
% org/location line" without ever looking like a spacing bug in the source.
% label=\textendash (not enumitem's default round bullet dot): the default
% bullet glyph has no working ToUnicode mapping under this lualatex+lmodern
% setup and extracts via pdftotext as a bare "�" replacement character in
% front of every single bullet -- a real ATS text-layer defect, caught during
% the mandatory pdftotext check on this template (2026-08-03). \textendash is
% already proven safe here (used in the section-heading dash rule below) and
% extracts cleanly as a literal "-".
\setlist[itemize]{leftmargin=1.15em, itemsep=0pt, topsep=0pt, parsep=0pt, partopsep=0pt, label=\textendash}

% Section headings: centered bold caps label flanked by a continuous rule on
% both sides, rather than a left-aligned label over a full-width rule. Built
% by hand instead of via titlesec. Two bugs found and fixed here 2026-08-03,
% surfaced by testing this design against a fully content-filled resume
% (cv/main_example_new_format_test.tex), not visible against the sparse
% placeholder text in this template alone:
% (1) v1 used a repeating dash-box pattern (\leaders\hbox to Npt{...}\hfil).
%     A fixed-width repeat can't guarantee it divides evenly into the
%     available stretch, so it often left a partial/clipped box at one
%     boundary -- fixed by using a single unbroken \hrule instead (no seams
%     to clip).
% (2) v1 built the heading as ordinary paragraph text (\noindent...\par).
%     That line is both the first AND last line of its own one-line
%     paragraph, so LaTeX's automatic \parfillskip (invisible "0pt plus
%     1fil" glue always appended at a paragraph's end) became a *third*
%     fil-order stretch competing with the two rule regions for the same
%     budget -- it skewed the split unevenly (one side visibly shorter than
%     the other) instead of 50/50, even though both rule regions had
%     identical stretch specs. Wrapping the whole line in
%     \makebox[\linewidth]{...} sidesteps paragraph-breaking (and
%     \parfillskip) entirely -- inside an explicit \hbox to <dim>, only the
%     glue actually present in that box competes for stretch, so the two
%     \hrulefill regions split evenly every time.
% CAUTION: \hrulefill has no shrinkability. A heading long enough to
% approach \linewidth on its own leaves the rule regions no room to work
% with and can overflow the margin (overfull hbox) instead of gracefully
% compressing -- keep section titles short, single-line phrases.
\newcommand{\rsection}[1]{%
  \par\vspace{9pt}%
  \noindent\makebox[\linewidth]{{\color{cvblue}\hrulefill}\enspace{\bfseries\color{cvblue}\MakeUppercase{#1}}\enspace{\color{cvblue}\hrulefill}}\par
  \vspace{4pt}%
}

% \resumeentry{Title/Role}{Date range}{Organization}{Location}
% Line 1: bold title, right-aligned bold date on the same line.
% Line 2: italic "Organization -- Location" (Location may be empty).
% \strut normalizes the line's height/depth regardless of which glyphs it
% contains, and \par (not a forced \\) ends the paragraph cleanly before the
% following itemize begins.
\newcommand{\resumeentry}[4]{%
  \noindent\textbf{#1}\hfill\textbf{#2}\\
  \ifstrempty{#4}{\textit{#3}}{\textit{#3 -- #4}}\strut\par
}

% personal data
% NOTE: never use an underscore in a bracketed placeholder token (e.g.
% [YOUR_NAME]) -- LaTeX treats a bare _ outside math mode as a subscript
% operator and throws "Missing $ inserted." Use a space instead: [YOUR NAME].
% A real, held credential can be appended directly, e.g. "[YOUR NAME], Ph.D."
% -- never append a credential the candidate doesn't actually hold.
\newcommand{\cvname}{[YOUR NAME]}
% Optional one-line headline under the name, e.g. your LinkedIn headline or
% target-title framing -- must be a real, verifiable line from the
% candidate's own profile, never an invented tagline. Delete the \cvheadline
% line in the HEADER block below if you don't want this row at all.
% Separator is a plain pipe (matching the contact line below), not a bullet
% dot -- $\bullet$ extracts as a bare "�" via pdftotext under this lualatex
% setup, the same replacement-character defect as the itemize bullet marker
% above (see that note for detail).
\newcommand{\cvheadline}{[TARGET TITLE] \textbar\ [KEY SPECIALTY / INDUSTRY]}
\newcommand{\cvphone}{[YOUR PHONE]}
\newcommand{\cvemail}{[YOUR EMAIL]}
\newcommand{\cvlinkedin}{[YOUR LINKEDIN URL]}
\newcommand{\cvaddress}{[YOUR ADDRESS]}

\begin{document}

% ============================================================
%     TOP ACCENT BAR CLEARANCE
% ============================================================
% The bar itself is drawn by the \AddToShipoutPictureBG hook in the preamble
% (a page-background overlay, independent of normal text flow) -- it does
% NOT push the rest of the content down automatically, so this \vspace does
% that job manually: clears the bar's 5pt height plus a bit of breathing
% room before the name.
% MUST be \vspace* (starred), not plain \vspace: ordinary \vspace glue
% placed at the very top of a page is silently discarded by TeX's
% page-breaking algorithm (a standard, easy-to-miss LaTeX quirk -- it's the
% same reason \topskip exists as its own mechanism rather than everyone just
% using \vspace). With plain \vspace here, the name would render directly
% under the bar, overlapping it -- caught only by comparing the compiled PDF
% against the reference image, not from the source looking wrong. \vspace*
% is specifically exempted from that discard rule.
\vspace*{20pt}

% ============================================================
%     HEADER
% ============================================================
% Name at \huge (up from \Large, 2026-08-03 per user request for a bigger
% name) -- the extra top clearance above (20pt, up from 15pt) accounts for
% \huge's taller cap-height so the name doesn't crowd the bar above it.
\begin{center}
{\huge\bfseries\color{cvblue}\cvname}\\[3pt]
{\small\cvheadline}\\[3pt]
{\small
% Plain-text contact line, no fontawesome icons -- fontawesome glyphs embed
% ToUnicode entries that make pdftotext/ATS extraction emit the glyph's
% internal name ("Envelope", "Map-mark") inline with the real contact text,
% which some ATS platforms surface directly to recruiters as parsed preview
% text. Confirmed via pdftotext -layout on a compiled resume.
\cvphone
\quad\textbar\quad
\href{mailto:\cvemail}{\cvemail}
\quad\textbar\quad
\href{\cvlinkedin}{LinkedIn}
\quad\textbar\quad
\cvaddress
}
\end{center}
\vspace{-8pt}

% ============================================================
%     PROFILE STATEMENT
% ============================================================
\needspace{4\baselineskip}
\rsection{Profile Statement}
\small{[1-3 sentences: who you are, your target job title verbatim from the posting, and your top qualifications for this specific role.]}

% ============================================================
%     CORE COMPETENCIES
% ============================================================
\needspace{5\baselineskip}
\rsection{Core Competencies}
\begin{itemize}
\item \textbf{[Competency Category 1]}: [comma-separated skills/tools]
\item \textbf{[Competency Category 2]}: [comma-separated skills/tools]
\item \textbf{[Competency Category 3]}: [comma-separated skills/tools]
\end{itemize}

% ============================================================
%     WORK EXPERIENCE
% ============================================================
\needspace{7\baselineskip}
\rsection{Work Experience}

\needspace{6\baselineskip}
\resumeentry{[Job Title]}{[Month YYYY -- Month YYYY]}{[Organization]}{[City, State]}
\begin{itemize}
\item \textbf{[2-4 word result label]:} Bullet emphasizing an achievement relevant to the target role, with a measurable result where possible -- the bold lead-in is optional but scans well; use only when it's a genuine, non-redundant label for the sentence that follows
\item Bullet
\item Bullet
\end{itemize}
\vspace{2pt}

\needspace{6\baselineskip}
\resumeentry{[Job Title]}{[Month YYYY -- Month YYYY]}{[Organization]}{[City, State]}
\begin{itemize}
\item Bullet
\item Bullet
\end{itemize}

% ============================================================
%     EDUCATION
% ============================================================
\needspace{5\baselineskip}
\rsection{Education}

\needspace{3\baselineskip}
\resumeentry{[Degree]}{[Month YYYY -- Month YYYY]}{[Institution]}{[City, State]}

\needspace{3\baselineskip}
\resumeentry{[Degree]}{[Month YYYY -- Month YYYY]}{[Institution]}{[City, State]}

% ============================================================
%     ADDITIONAL SECTIONS (optional -- add only if the 1-page budget allows:
%     Languages, Publications, Honors and Awards. Do NOT add a References
%     section -- omit it entirely at this length; the standalone reference
%     sheet (08-reference-sheet-templates.md) covers this for every
%     application regardless of what the resume itself says.)
% ============================================================

\end{document}
